\documentclass[11pt]{article}

\usepackage{nmrdoc}
\usepackage{siunitx}
\usepackage{bm}

\sisetup{per-mode=symbol}
\DeclareSIUnit{\angstrom}{\text{\AA}}

\begin{document}
\NmrTitle{The Haigh--Mallion Ring-Current Surface Calculator}

\section{What the calculator computes}

The Haigh--Mallion calculator represents the magnetic shielding pattern caused
by an aromatic ring current. When an external field is applied to an aromatic
ring, the delocalized \(\pi\) electrons can circulate around the ring. That
circulation produces a secondary magnetic field at a nearby nucleus. The
secondary field changes the nuclear shielding tensor \(\bm\sigma\), and the
isotropic part of \(\bm\sigma\) contributes to the chemical shift reported in
solution NMR.

This implementation does not solve the electronic response of the ring. It
takes the ring geometry from the protein structure, integrates a magnetic
dipole kernel over the polygonal ring surface, and stores the resulting
geometric tensor. The output has units \(\si{\angstrom^{-1}}\). It is not a
reported shielding value in ppm and does not stand on its own as a chemical
shift. After calibration against DFT shielding tensors, it can contribute a
ring-current shielding term; before calibration, it is a geometry-derived
descriptor whose tensor shape is the object of interest.

\section{Where shielding enters}

For a nucleus in an applied magnetic field \(\bm B_0\), the field experienced
by the nucleus is written
\begin{equation}
  \bm B_{\mathrm{nuc}}=(\bm I-\bm\sigma)\bm B_0 .
  \label{eq:shielding}
\end{equation}
Here \(\bm I\) is the \(3\times3\) identity matrix. The shielding tensor
\(\bm\sigma\) is dimensionless in physical NMR: it is the linear-response
coefficient connecting the applied field to the electronic secondary field at
the nucleus. It is a matrix because the molecule has orientation. A field
applied along direction \(b\) can induce a secondary-field component along
direction \(a\), so \(\sigma_{ab}\) records how those two directions are
coupled.

The quantum-mechanical content behind Eq.~\eqref{eq:shielding} is that the
external magnetic field perturbs the electronic state. The perturbed electrons
generate a magnetic field of their own, and the nucleus responds to the sum.
The Haigh--Mallion calculator keeps one classical, geometry-based piece of that
response. For an aromatic ring, the component of \(\bm B_0\) normal to the ring
threads magnetic flux through the \(\pi\)-electron system. The model represents
the induced current response by a thin magnetic surface lying inside the ring.

The sign convention in the code is
\begin{equation}
  \sigma_{ab}=-\frac{\partial B^{\mathrm{sec}}_a}{\partial B_{0,b}} .
  \label{eq:sign}
\end{equation}
The minus sign says that shielding is the reduction of the applied field at the
nucleus. In the source, this sign appears when the full kernel \(G_{ab}\) is
formed from the surface integral. A diamagnetic literature ring-current
intensity is negative, but that intensity is not multiplied into this
calculator's stored result.

\section{The tensor split used by the code}

Any \(3\times3\) tensor \(\bm X\) is decomposed by
\texttt{SphericalTensor::Decompose} before it is emitted. The scalar part is
\begin{equation}
  T_0=\frac{1}{3}\operatorname{tr}\bm X .
  \label{eq:t0}
\end{equation}
This is the isotropic component: it multiplies \(\bm I\), survives molecular
tumbling, and sets the isotropic shielding value before reference subtraction.

The antisymmetric part is packed as the three-component vector
\begin{equation}
  \bm T_1=\frac{1}{2}
  \begin{pmatrix}
    X_{yz}-X_{zy}\\
    X_{zx}-X_{xz}\\
    X_{xy}-X_{yx}
  \end{pmatrix}.
  \label{eq:t1}
\end{equation}
This vector is the dual form of the part of \(\bm X\) that changes sign under
transpose. It is usually not observed directly in standard isotropic solution
NMR, but it is retained because the Haigh--Mallion shielding kernel need not be
symmetric.

The remaining symmetric traceless matrix is
\begin{equation}
  \bm S=\frac{\bm X+\bm X^{\mathsf T}}{2}-T_0\bm I,
  \qquad \operatorname{tr}\bm S=0 .
  \label{eq:symmetric-traceless}
\end{equation}
It has five independent entries. The code stores those entries in the
real-spherical \(T_2\) basis
\begin{equation}
  \bm T_2=
  \begin{pmatrix}
    \sqrt{2}\,S_{xy}\\
    \sqrt{2}\,S_{yz}\\
    \sqrt{3/2}\,S_{zz}\\
    \sqrt{2}\,S_{xz}\\
    (S_{xx}-S_{yy})/\sqrt{2}
  \end{pmatrix}.
  \label{eq:t2}
\end{equation}
These factors are the normalization used by the implementation. They preserve
the Frobenius norm, so the squared length of the five-vector equals
\(\sum_{ab}S_{ab}^2\). The nine stored values are ordered as
\[
  [T_0,T_{1x},T_{1y},T_{1z},
    T_{2,-2},T_{2,-1},T_{2,0},T_{2,+1},T_{2,+2}] .
\]

The Haigh--Mallion path creates two tensors. The raw surface integral
\(\bm H\) is symmetric and traceless up to floating-point quadrature error, so
its decomposition is carried by \(T_2\). The stored shielding kernel
\(\bm G\), formed after contraction with the ring normal, is a rank-one matrix
and can have nonzero \(T_0\), \(T_1\), and \(T_2\). The per-atom output uses the
decomposition of the accumulated \(\bm G\), not the raw \(\bm H\).

\section{The method implemented}

For each protein conformation, the calculator depends on the precomputed ring
geometry and a spatial index. The ring list is the aromatic part of the protein
topology: PHE benzene, TYR phenol, TRP benzene, TRP pyrrole, TRP indole
perimeter, HIS imidazole, HID imidazole, and HIE imidazole. A saturated PRO
pyrrolidine type exists elsewhere in the topology, but the aromatic ring-count
API used here does not present it to this loop.

For a ring with vertices \(\bm q_i\), the geometry routine sets the center to
the average vertex position,
\begin{equation}
  \bm c=\frac{1}{N}\sum_{i=1}^{N}\bm q_i ,
  \label{eq:center}
\end{equation}
with positions measured in \(\si{\angstrom}\). The ring normal
\(\hat{\bm n}\) is the normal to the best-fit plane from an SVD of the centered
vertex coordinates. Its sign is chosen to agree with
\((\bm q_1-\bm q_0)\times(\bm q_2-\bm q_0)\). The radius used by filters is
the average distance of the vertices from \(\bm c\).

For a target atom at position \(\bm r\), the raw tensor is
\begin{equation}
  H_{ab}(\bm r)
  =\int_{S}
  \left(
    \frac{3\rho_a\rho_b}{\rho^5}
    -\frac{\delta_{ab}}{\rho^3}
  \right)\,dS,
  \qquad
  \bm\rho=\bm r-\bm s,\quad \rho=|\bm\rho| .
  \label{eq:h-integral}
\end{equation}
The integration point \(\bm s\) runs over the ring surface \(S\);
\(dS\) is an area element in \(\si{\angstrom^2}\);
\(\delta_{ab}\) is one when \(a=b\) and zero otherwise. The integrand has units
\(\si{\angstrom^{-3}}\), so \(\bm H\) has units \(\si{\angstrom^{-1}}\).
The matrix inside the integral is the Cartesian magnetic-dipole field kernel:
if a unit point dipole at \(\bm s\) points along direction \(b\), the expression
gives the field component along direction \(a\) at \(\bm r\), apart from
physical constants not stored here. The Haigh--Mallion surface is therefore a
sheet of infinitesimal dipoles. Integrating Eq.~\eqref{eq:h-integral} asks what
field that sheet would produce at the nucleus for each possible dipole
direction.

The ring current is tied to the ring normal. The implementation contracts
\(\bm H\) with \(\hat{\bm n}\),
\begin{equation}
  \bm V(\bm r)=\bm H(\bm r)\hat{\bm n},
  \label{eq:v-field}
\end{equation}
so \(V_a\) is the effective secondary-field component from a unit surface
response normal to the ring. The applied field enters through its normal
component, \(\hat{\bm n}\cdot\bm B_0=n_bB_{0,b}\). Combining that flux factor
with the shielding sign convention gives the stored tensor
\begin{equation}
  G_{ab}(\bm r)=-V_a(\bm r)n_b
  =-\left[\bm H(\bm r)\hat{\bm n}\right]_a n_b .
  \label{eq:g-kernel}
\end{equation}
This is an outer product, so the Cartesian matrix has rank one unless
\(\bm V=\bm 0\). Its units remain \(\si{\angstrom^{-1}}\), because
\(\hat{\bm n}\) is dimensionless.

The surface \(S\) is not meshed by an external library. The code creates a fan
triangulation from the ring center to each adjacent vertex pair. On each
triangle \(T=(\bm v_0,\bm v_1,\bm v_2)\), it evaluates
\begin{equation}
  \int_T F(\bm s)\,dS
  \approx
  A_T\sum_{q=1}^{7}w_q
  F(\lambda_{q0}\bm v_0+\lambda_{q1}\bm v_1+\lambda_{q2}\bm v_2),
  \label{eq:quadrature}
\end{equation}
where \(A_T\) is the triangle area. The \(\lambda\)'s are barycentric
coordinates: nonnegative weights that sum to one and locate a point inside the
triangle. The seven quadrature weights also sum to one. The rule is the
degree-five Stroud T2:5-1, also listed as Dunavant's seven-point triangle rule;
the code computes its two barycentric orbit parameters
\((6\pm\sqrt{15})/21\) and weights \(9/40\) and
\((155\pm\sqrt{15})/1200\).

Near the ring plane, the \(\rho^{-3}\) kernel changes rapidly. The
implementation refines a fan triangle by splitting it at edge midpoints when
the target atom is close to a triangle vertex. The first split is triggered
below \(\SI{2.0}{\angstrom}\), the second below \(\SI{1.0}{\angstrom}\), and
the depth stops at two. That leaves 7 quadrature points without subdivision,
28 after the first split, and at most 112 if all level-1 subtriangles split. A
triangle whose area is below \(10^{-20}\,\si{\angstrom^2}\) is skipped, and a
quadrature point closer than \(\SI{0.1}{\angstrom}\) to the target atom is
skipped.

Candidate ring centers are first found within \(\SI{15.0}{\angstrom}\) of the
target atom. Three filters then decide whether the atom--ring pair is
evaluated. The \(\SI{0.1}{\angstrom}\) singularity guard rejects a target
closer than that distance to the expansion center. A near-field filter rejects
points whose center distance is not larger than half the source extent; for
rings the source extent is the diameter, so the threshold is the ring radius. A
topology filter rejects ring vertices and atoms covalently bonded to ring
vertices, because those atoms lie in a through-bond regime rather than the
through-space surface model.

The TOML file also contains ring-current intensities in \(\si{\nano\ampere\per
\tesla}\): PHE \(-12.0\), TYR \(-11.28\), TRP6 \(-12.48\), TRP5 \(-6.72\),
HIS/HID/HIE \(-5.16\), and TRP9 \(-19.2\). It contains Johnson--Bovey lobe
offsets in \(\si{\angstrom}\): \(0.64\) for PHE, TYR, and TRP6; \(0.52\) for
TRP5; \(0.50\) for HIS/HID/HIE; and \(0.60\) for TRP9. This calculator reads
neither set. The source computes \(\bm G\) from geometry, records the ring type
for per-type sums, and lets calibration or a later model decide how strongly
each ring type should contribute.

\section{Inputs and output}

The direct inputs are atomic coordinates, typed aromatic-ring topology, the
ring geometry derived from those coordinates, and the spatial index used to
find nearby rings. The calculator does not consume MOPAC charges or bond
orders, AIMNet2 charges or embeddings, ff14SB charges, APBS electrostatic
fields, or DSSP secondary-structure assignments. Those data sources feed other
calculators or provide structural context; their internal methods are outside
the scope of this document.

For each atom, the calculator stores per-ring attribution: \(\bm H\), its
decomposition, \(\bm V=\bm H\hat{\bm n}\), \(\bm G\), and the decomposition of
\(\bm G\). It accumulates the accepted \(\bm G\) tensors over nearby rings and
writes \texttt{hm\_shielding.npy} with shape \(N\times9\), using the
\([T_0,T_1,T_2]\) order given after Eq.~\eqref{eq:t2}. It also writes
\texttt{hm\_per\_type\_T0.npy} with shape \(N\times8\) and
\texttt{hm\_per\_type\_T2.npy} with shape \(N\times40\), one \(T_0\) scalar
and five \(T_2\) components for each aromatic ring type.

The grid-sampling method evaluates the same \(\bm G\) construction at an
arbitrary point. It does not use the bonded topology filter, and it skips rings
closer than the singularity guard, points inside the ring radius, and rings
beyond \(\SI{15.0}{\angstrom}\).

Despite the output filename, the values are geometric kernels. Their role in
the thesis is to test whether the Haigh--Mallion geometry carries the tensor
signal seen in DFT shielding differences. Agreement is therefore assessed as
correlation after calibration, especially in the \(T_2\) anisotropic
components, not as pointwise reproduction of a chemical shift.

\end{document}
